% !PRINTED 95
\weberpage{95}
If, then, we understand by $q_1, q_2, q_3 \ldots$ numbers taken from the sequence consisting of the number 9 and the primes 7, 13, 19, 31 . . . , then the general expression for $m$ is:
\weq{1}{m = q_1\,q_2\,q_3\cdots q_\mu.}
The index moduli are, when $q_1=9$, $c_1=6$, $c_2=q_2-1$, $c_3=q_3-1 \ldots$ If 9 and 7 do not occur among the factors of $m$, the index modulus 6 drops out.

The quantities $\delta_{1,1}, \delta_{1,2}, \ldots \delta_{1,\mu}$ are the greatest common divisors of $i_1=3$ with the $c_1, c_2, \ldots c_\mu$; they are therefore all $=3$, and by \S.\,21, (11) it follows that:
\[
i_{1,h}=1,\quad c_{1,h}=\tfrac{1}{3}\,c_1.
\]
% [sic: Weber prints c_1 on the right; because h is arbitrary and c_{1,h}=c_h/delta_{1,h}, the intended subscript is h.]
The matrix \S.\,21, (24) here consists of only the single vertical row $e_{1,1}, e_{1,2}, \ldots e_{1,\mu}$, and these quantities may be chosen arbitrarily, provided only that at least one is among them that is indivisible by 3. These numbers $e_{1,1}, e_{1,2}, \ldots e_{1,\mu}$ are, moreover, determined only to the modulus 3, and may therefore be taken equal to 0, 1 or 2. But condition 6. of \S.\,21 also requires that, when $q_1=9$, $c_{1,1}e_{1,1}=2e_{1,1}$ shall not be divisible by 3, and, when $q_2$ is a prime, that
\[
c_{1,2}\,e_{1,2}=\tfrac{1}{3}(q_2-1)\,e_{1,2}
\]
shall not be divisible by $q_2-1$; that is to say, \gesperrt{none of the numbers} $e_{1,1}, e_{1,2}, \ldots e_{1,\mu}$ may be divisible by 3, and we can choose $\pm 1$ for each of them.

In order, then, to find the indices $\alpha_1, \alpha_2, \ldots \alpha_\mu$ of all the numbers $a$ that can occur in the period
\weq{2}{\eta = \sum r^a}
one has, by \S.\,22, (3), to take all the numbers $\alpha$ (to the modulus 3) that satisfy a congruence
\weq{3}{\pm\alpha_1\pm\alpha_2\pm\alpha_3\cdots\pm\alpha_\mu\equiv 0\ (\mathrm{mod}\ 3)}
where for \gesperrt{a definite group one combination of signs must be held fixed}.

The number of the possible cubic fields depends, for given $m$, only on the number $\mu$ of the index moduli, and, since one of the signs in (3) can be assumed arbitrarily, amounts to $2^{\mu-1}$.

The result would take an entirely similar shape if one were to put another prime $p$ in place of 3 for $i_1$. There would then merely enter into the congruence (3), in place of the double
%CONT>>
